% !TeX program = pdflatex
% Tom III: Analiza stochastyczna; rozdzialy 1--16, audyt rozdzialow 13--14.
% Przygotowano jako samodzielny dokument; zasady i symbole jak w tomie II.
\documentclass[11pt,a4paper,oneside]{report}
\usepackage[T1]{fontenc}
\usepackage[utf8]{inputenc}
\usepackage[polish]{babel}
\usepackage{lmodern}
\usepackage{amsmath,amssymb,amsthm,mathtools}
\usepackage{microtype}
\usepackage{geometry}
\usepackage{enumitem}
\usepackage{xcolor}
\usepackage{tikz}
\usetikzlibrary{arrows.meta,positioning,calc,shapes.geometric,decorations.pathmorphing}
\usepackage[unicode]{hyperref}
% Odsylacze miedzytomowe; zbudowac wszystkie .aux i przechowywac PDF-y obok siebie.
\usepackage{xr-hyper}
\externaldocument[I-]{rachunek_prawdopodobienstwa_tom_I}[rachunek_prawdopodobienstwa_tom_I.pdf]
\externaldocument[II-]{rachunek_prawdopodobienstwa_tom_II}[rachunek_prawdopodobienstwa_tom_II.pdf]
\externaldocument[IV-]{rachunek_prawdopodobienstwa_tom_IV}[rachunek_prawdopodobienstwa_tom_IV.pdf]

\geometry{margin=2.65cm}
\hypersetup{colorlinks=true,linkcolor=blue!45!black,urlcolor=blue!45!black}
\setlength{\parindent}{0pt}
\setlength{\parskip}{4pt}
\setlength{\emergencystretch}{2em}
\setlist[enumerate]{topsep=3pt,itemsep=2pt}
\setlist[itemize]{topsep=3pt,itemsep=2pt}
\numberwithin{equation}{chapter}
\theoremstyle{plain}
\newtheorem{twierdzenie}{Twierdzenie}[chapter]
\newtheorem{lemat}[twierdzenie]{Lemat}
\newtheorem{wniosek}[twierdzenie]{Wniosek}
\newtheorem{stwierdzenie}[twierdzenie]{Stwierdzenie}
\theoremstyle{definition}
\newtheorem{definicja}{Definicja}[chapter]
\newtheorem{przyklad}{Przykład}[chapter]
\newtheorem{zadanie}{Zadanie rozwiązane}[chapter]
\theoremstyle{remark}
\newtheorem{uwaga}{Uwaga}[chapter]
\renewcommand{\proofname}{Dowód}
\newcommand{\R}{\mathbb{R}}
\newcommand{\N}{\mathbb{N}}
\newcommand{\Z}{\mathbb{Z}}
\newcommand{\Q}{\mathbb{Q}}
\newcommand{\Prob}{\mathbb{P}}
\newcommand{\E}{\mathbb{E}}
\newcommand{\Var}{\operatorname{Var}}
\newcommand{\Cov}{\operatorname{Cov}}
\newcommand{\Law}{\operatorname{Law}}
\newcommand{\calF}{\mathcal{F}}
\newcommand{\calE}{\mathcal{E}}
\newcommand{\calH}{\mathcal{H}}
\newcommand{\calN}{\mathcal{N}}
\newcommand{\calB}{\mathcal{B}}
\newcommand{\calP}{\mathcal{P}}
\newcommand{\calG}{\mathcal{G}}
\newcommand{\calD}{\mathcal{D}}
\newcommand{\calS}{\mathcal{S}}
\newcommand{\calC}{\mathcal{C}}
\newcommand{\ind}[1]{\mathbf{1}_{#1}}
\newcommand{\dd}{\mathrm{d}}
\newcommand{\norm}[1]{\lVert #1\rVert}
\newcommand{\abs}[1]{\lvert #1\rvert}
\newcommand{\iid}{\mathrel{\overset{\mathrm{iid}}{\sim}}}
\newcommand{\Normal}{\mathcal{N}}
\DeclareMathOperator{\dist}{dist}
\DeclareMathOperator{\supp}{supp}
\DeclareMathOperator{\spanop}{span}